\documentclass[11pt,a4paper]{article}

% ----------------------------------------------------------- Pakete
\usepackage[ngerman]{babel}
\usepackage[T1]{fontenc}
\usepackage[utf8]{inputenc}
\usepackage{lmodern}
\usepackage{microtype}
\usepackage{xcolor}
\usepackage{geometry}
\geometry{margin=2.5cm, headsep=12pt}
\usepackage{amsmath,amssymb,mathtools}
\usepackage{graphicx}
\usepackage{booktabs}
\usepackage{siunitx}
\usepackage[font=small,labelfont=bf]{caption}
\usepackage{subcaption}
\usepackage{enumitem}
\usepackage{listings}
\usepackage{tikz}
\usetikzlibrary{arrows.meta,positioning,shapes.geometric}
\usepackage{pgfplots}
\pgfplotsset{compat=1.18}
\usepackage{fancyhdr}
\usepackage{hyperref}
\hypersetup{
  colorlinks=true,
  linkcolor=blue!45!black,
  citecolor=green!40!black,
  urlcolor=blue!60!black,
  pdftitle={Messung von Satzqualität},
}
\usepackage{cleveref}

\lstset{
  basicstyle=\ttfamily\small,
  breaklines=true,
  frame=single,
  rulecolor=\color{black!30},
  keywordstyle=\color{blue!60!black}\bfseries,
  commentstyle=\color{green!40!black}\itshape,
  showstringspaces=false,
  language=[LaTeX]TeX,
}

\pagestyle{fancy}
\fancyhf{}
\fancyhead[L]{\small Satzqualität -- Demoartikel}
\fancyhead[R]{\small \thepage}
\renewcommand{\headrulewidth}{0.4pt}

\title{Messung von Satzqualität\\[0.4em]\large Ein vollständiges \LaTeX-Demoprojekt}
\author{Erika Beispiel \and Max Muster}
\date{\today}

\begin{document}
\maketitle

\begin{abstract}
  Dieses Demoprojekt zeigt alle Funktionen von \emph{LaTeX Studio} an einem
  realistischen Anwendungsfall: mehrere Abbildungen, Tabellen, Programmcode,
  Verweise auf Gleichungen und Abschnitte, Zitate aus einer Literaturdatenbank
  sowie eine automatisch erzeugte Gliederung. Es eignet sich hervorragend zum
  Ausprobieren von Auto-Build, SyncTeX und der Projektweiten Suche.
\end{abstract}

\begin{figure}[htbp]
  \centering
  \begin{tikzpicture}[
      node distance=16mm and 20mm,
      box/.style={draw, rounded corners=2pt, fill=blue!8, minimum height=9mm, align=center},
      arr/.style={-{Stealth[length=2.6mm]}, thick, black!60},
    ]
    \node[box] (src) {Quelltext\\(.tex)};
    \node[box, right=of src] (eng) {Engine\\(pdf\LaTeX)};
    \node[box, right=of eng] (pdf) {PDF};
    \node[box, below=of pdf, fill=orange!12] (sync) {SyncTeX};
    \node[box, below=of src, fill=green!10] (ui) {Editor};
    \draw[arr] (src) -- (eng);
    \draw[arr] (eng) -- (pdf);
    \draw[arr] (pdf) -- (sync);
    \draw[arr] (sync) -- (ui);
    \draw[arr] (ui) -- (src);
  \end{tikzpicture}
  \caption{Der Build-Kreislauf von \emph{LaTeX Studio} mit SyncTeX-Rücksprung.}
  \label{fig:pipeline}
\end{figure}

\section{Einleitung}
\label{sec:intro}

Satzqualität entsteht dort, wo Trennung von Inhalt und Gestaltung konsequent
durchgezogen wird~\cite{lamport1994}. Während Textverarbeitungen Formatierung
am Absatz speichern, beschreibt \LaTeX\ die \emph{Bedeutung} eines Bausteins;
die Gestaltung erfolgt in Klassen und Paketen~\cite{kopka2003,mittelbach2004}.

\cref{fig:pipeline} zeigt den Kreislauf aus Quelltext, Kompilierung und
Rücksprung in den Editor. In \cref{sec:methode} beschreiben wir die Messgrößen,
in \cref{sec:ergebnisse} folgen die Ergebnisse.

\section{Messverfahren}
\label{sec:methode}

\subsection{Metriken}

Wir bewerten ein Dokument mit dem gewichteten Summenmaß
\begin{equation}
  \label{eq:score}
  Q \;=\; \frac{\sum_{i=1}^{n} w_i \, s_i}{\sum_{i=1}^{n} w_i},
  \qquad w_i > 0,\; s_i \in [0,1],
\end{equation}
wobei $s_i$ den Wert der $i$-ten Metrik und $w_i$ deren Gewicht bezeichnet.
Eine ergänzende Fehlermaßzahl ist die mittlere Abweichung
\begin{align}
  \bar{e} &= \frac{1}{n}\sum_{i=1}^{n}\bigl|s_i - \hat{s}_i\bigr|
  \label{eq:mae}\\
  \operatorname{RMSE} &= \sqrt{\frac{1}{n}\sum_{i=1}^{n}\bigl(s_i - \hat{s}_i\bigr)^2}.
  \label{eq:rmse}
\end{align}
Nach \cref{eq:score} ist $Q$ genau dann groß, wenn alle Gewichte auf
hohe Einzelwerte zielen.

\subsection{Implementierung}

Die Auswertung ist als kurzes \LaTeX-Makro skizziert in
\cref{lst:makro}.

\begin{lstlisting}[caption={Gewichtete Mittelwertbildung}, label={lst:makro}]
\newcommand{\gewichtet}[1]{%
  \num{\frac{#1}{\n}}%
}
\end{lstlisting}

\section{Ergebnisse}
\label{sec:ergebnisse}

\subsection{Vergleich der Verfahren}

\begin{table}[htbp]
  \centering
  \caption{Ergebnisse der Qualitätsmessung (Mittelwerte über 30 Läufe).}
  \label{tab:ergebnisse}
  \begin{tabular}{l S[table-format=1.3] S[table-format=1.3] S[table-format=1.2]}
    \toprule
    Verfahren & {$Q$} & {$\bar{e}$} & {Laufzeit (\si{\second})} \\
    \midrule
    Regelbasiert   & 0,712 & 0,148 & 1,24 \\
    Lerndbasiert   & 0,847 & 0,091 & 3,86 \\
    Hybrid         & 0,871 & 0,082 & 3,41 \\
    \bottomrule
  \end{tabular}
\end{table}

\subsection{Verlauf über die Epochen}

\begin{figure}[htbp]
  \centering
  \begin{tikzpicture}
    \begin{axis}[
        width=0.86\linewidth,
        height=6cm,
        xlabel={Epoche},
        ylabel={Qualität $Q$},
        xmin=1, xmax=10,
        ymin=0.4, ymax=1.0,
        grid=major,
        grid style={black!12},
        legend style={at={(0.02,0.98)}, anchor=north west, font=\small, draw=black!30},
      ]
      \addplot[blue, thick, mark=*, mark size=1.8] coordinates {
        (1,0.52) (2,0.61) (3,0.68) (4,0.73) (5,0.77)
        (6,0.80) (7,0.83) (8,0.85) (9,0.86) (10,0.871)
      };
      \addplot[red, thick, dashed, mark=square*, mark size=1.6] coordinates {
        (1,0.50) (2,0.56) (3,0.62) (4,0.66) (5,0.70)
        (6,0.72) (7,0.74) (8,0.75) (9,0.76) (10,0.768)
      };
      \legend{Hybrid, Regelbasiert}
    \end{axis}
  \end{tikzpicture}
  \caption{Konvergenz über zehn Epochen. Die Kurven entsprechen den Werten aus
    \cref{tab:ergebnisse}.}
  \label{fig:konvergenz}
\end{figure}

Wie \cref{fig:konvergenz} zeigt, nähern sich beide Verfahren einem Wert von
$Q \approx 0{,}87$ an; der hybride Ansatz liegt dabei durchgehend oben.

\section{Diskussion}
\label{sec:diskussion}

\begin{itemize}[leftmargin=*]
  \item Die Gewichtung in \cref{eq:score} empfindlich zu wählen, verändert das
    Ranking der Verfahren.
  \item \cref{eq:rmse} bestraft Ausreißer stärker als \cref{eq:mae}.
  \item Die Laufzeiten in \cref{tab:ergebnisse} skalieren linear mit der Zahl
    der Metriken.
\end{itemize}

Offen bleibt die Frage, ob sich die Gewichte $w_i$ automatisch aus
Bestandsdaten ableten lassen~\cite{mittelbach2004}.

\section{Fazit}
\label{sec:fazit}

Mit einem einzigen Quelltext entstehen Gliederung, Abbildungen, Tabellen,
Code listings und ein konsistentes Literaturverzeichnis. Genau das macht
\emph{LaTeX Studio} zum passenden Werkzeug für dieses Vorgehen.

% TODO: Abschnitt zu Ausblick und Reproduzierbarkeit ergänzen

\clearpage
\bibliographystyle{plain}
\bibliography{references}

\end{document}
